% ---------------------------------------------------------------------------
% Résumé en français et mots-clés
% ---------------------------------------------------------------------------

\chapter*{Résumé}

Le produit Scan d'ISOGEO permet d'explorer différentes sources de données afin
d'identifier les ressources qu'elles contiennent et d'en extraire les
informations nécessaires à leur documentation.
Son fonctionnement repose sur
plusieurs traitements capables d'analyser des bases de données, des
géodatabases et des fichiers géographiques provenant d'environnements
techniques variés.

Une partie de ces traitements était historiquement réalisée à l'aide de
workbenches FME. Bien que cet outil fournisse de nombreuses fonctionnalités
d'intégration et de transformation de données, son utilisation introduisait une
dépendance externe dans la chaîne d'exécution de Scan.
La mission réalisée au
cours de cette alternance a donc consisté à remplacer progressivement certains
de ces traitements par une solution Python intégrée au produit et pouvant être
exécutée de manière autonome.

Le travail a d'abord porté sur l'analyse des traitements FME existants afin
d'identifier leurs entrées, leurs sorties, les transformations appliquées et les
règles métier à conserver.
Des composants Python ont ensuite été développés
pour se connecter à différentes sources, notamment PostgreSQL/PostGIS, Oracle,
SQL Server et les géodatabases d'entreprise Esri.
Les développements réalisés
permettent d'extraire des informations relatives aux tables, aux champs, aux
projections, aux emprises spatiales, aux géométries et à différentes
métadonnées techniques attendues par Scan.

La bibliothèque GDAL/OGR a notamment été utilisée pour l'analyse des sources
géographiques vectorielles et raster.
Des traitements spécifiques ont également
été mis en place pour gérer les particularités de chaque système de gestion de
base de données, les données invalides, les erreurs de connexion et les écarts
de comportement entre les différents formats.

Une attention particulière a été portée à l'industrialisation de la solution.
Les composants développés ont été intégrés dans une architecture modulaire,
accompagnés d'une gestion structurée des erreurs et préparés pour la génération
d'un exécutable Windows autonome avec PyInstaller.
Une démarche de tests de
régression a également été élaborée afin de comparer automatiquement les
résultats produits par les anciens traitements FME avec ceux obtenus par la
nouvelle implémentation Python.

Les travaux réalisés contribuent ainsi à réduire la dépendance de Scan à FME,
à faciliter la maintenance de ses traitements et à renforcer leur intégration
dans l'environnement logiciel d'ISOGEO. Cette alternance m'a également permis
de développer mes compétences en programmation Python, en traitement de
données géographiques, en bases de données spatiales, en packaging logiciel et
en validation de solutions techniques.

\vspace{0.7cm}

\noindent
\textbf{Mots-clés :}
géomatique ; Scan ; ISOGEO ; Python ; FME ; GDAL/OGR ; métadonnées ;
bases de données spatiales ; PyInstaller ; tests de régression.
